\documentclass[reqno,a4paper,12pt]{amsart}
\usepackage[OT1]{fontenc}
\usepackage{amsmath,amssymb,amsthm}
\usepackage[textwidth=152mm,textheight=230mm,centering]{geometry}
\usepackage{microtype}
\usepackage{xurl}
\usepackage[hidelinks]{hyperref}
\hypersetup{pdftitle={Schauder bases in unitary orbits}}
\newtheorem{theorem}{Theorem}
\newtheorem{proposition}[theorem]{Proposition}
\newtheorem{lemma}[theorem]{Lemma}
\newcommand{\R}{\mathbb R}
\newcommand{\C}{\mathbb C}
\newcommand{\T}{\mathbb T}
\newcommand{\Q}{\mathbb Q}
\newcommand{\Z}{\mathbb Z}
\newcommand{\E}{\mathbb E}
\newcommand{\spn}{\operatorname{span}}
\newcommand{\dist}{\operatorname{dist}}
\newcommand{\diag}{\operatorname{diag}}
\newcommand{\Var}{\operatorname{Var}}
\newcommand{\ip}[2]{\langle #1,#2\rangle}
\newcommand{\norm}[1]{\left\lVert #1\right\rVert}
\newcommand{\ml}[1]{\nolinkurl{#1}}
\title{Schauder bases in unitary orbits}
\date{November 18, 2026}
\begin{document}
\begin{abstract}
A strongly continuous one-parameter unitary group on a separable
infinite-dimensional complex Hilbert space has simple continuous spectrum
if and only if one of its orbits contains a Schauder basis.
The specialized analytic and finite-dimensional constructions are proved
here; the Mathlib foundations used are identified in the dependency audit.
\end{abstract}
\maketitle

\section{Statement and conventions}
Let $U=(U_t)_{t\in\R}$ be a strongly continuous one-parameter group of
unitary operators on a separable infinite-dimensional complex Hilbert space $H$.
A vector $v$ is cyclic if $\spn\{U_tv:t\in\R\}$ is dense in $H$.
A common eigenvector is a nonzero $a$ such that $U_ta\in\C a$ for every $t$.
Simple continuous spectrum means the existence of a cyclic vector and
the absence of common eigenvectors.
A Schauder basis is a sequence with a unique norm-convergent scalar
expansion for every vector.

\begin{theorem}\label{thm:main}
The group $U$ has simple continuous spectrum if and only if there are
$g\in H$ and distinct $t_n\in\R$ such that $(U_{t_n}g)_{n\ge1}$
is a Schauder basis of $H$.
\end{theorem}

Inner products are linear in the first variable. Write
$e_t(x)=\exp(2\pi itx)$ and equip $\T^d=(\R/\Z)^d$ with normalized
Haar measure. All measures on $\R$ below are Borel measures. All weights
are measurable; equalities between $L^2$ functions are almost everywhere.
The only mathematical results used without proof beyond elementary
algebra, calculus and measure-space definitions are the Mathlib foundations
listed in Section~\ref{audit}. In particular, neither a spectral
representation theorem nor a weighted exponential-basis theorem is assumed.

\section{Bases, characters, and the spectral measure}

\begin{lemma}[The needed basis facts]\label{basisfacts}
For a Schauder basis $(b_n)$ of a Banach space, its coefficient maps are
continuous and its partial-sum maps are uniformly bounded.
In a Hilbert space a bounded Schauder basis is weakly null.
Conversely, a linearly independent sequence with dense span is a Schauder
basis if all its partial-sum maps on the algebraic span have a common bound.
\end{lemma}
\begin{proof}
Uniqueness implies $b_n\ne0$ and finite linear independence. Let $X$ be the
space of scalar sequences $a$ for which $\sum a_nb_n$ converges, with norm
$\|a\|_X=\sup_N\|\sum_{n\le N}a_nb_n\|$.
For every $n$, $|a_n|\|b_n\|\le2\|a\|_X$.
If $a^{(j)}$ is Cauchy in $X$, each coordinate has a limit $a_n$.
Passing to the limit in each finite sum gives
$\sup_N\|\sum_{n\le N}(a_n-a_n^{(j)})b_n\|\le\epsilon$
for all sufficiently large $j$. The partial sums for $a$ are Cauchy:
their tail differences are at most $2\epsilon$ plus the corresponding
tail difference for a fixed $a^{(j)}$. Thus $a\in X$ and
$a^{(j)}\to a$ in $X$. This proves completeness.
The sum map $X\to H$ is a continuous linear bijection, so its inverse
is continuous by the bounded inverse theorem. Consequently
$\sup_N\|P_Nx\|\le K\|x\|$ and each coefficient is continuous.

In the Hilbert case represent the coefficient functionals as
$b_n^*(x)=\ip{x}{y_n}$ by the Hilbert-space Riesz theorem.
If $x$ is orthogonal to every $y_n$, all its coefficients vanish, so $x=0$.
The span of the $y_n$ is therefore dense. For $M=\sup_n\|b_n\|<\infty$
and $y$ within $\epsilon$ of $\spn\{y_1,\ldots,y_N\}$,
$|\ip{b_n}{y}|\le M\epsilon$ for $n>N$ by biorthogonality.
This is weak convergence to zero, without using a shrinking-basis theorem.

For the converse, if $p_j$ in the algebraic span tend to $x$, define
$P_Nx=\lim_jP_Np_j$. The bound $K$ makes this a Cauchy sequence and makes
its limit independent of the approximating sequence. Linearity and the
same bound pass to the limit.
Approximation by a finite sum $p$ gives
$\|P_Nx-x\|\le(K+1)\|x-p\|$ for large $N$.
The differences $P_n-P_{n-1}$ have range $\C b_n$, hence determine
continuous scalar coefficients. Their sums converge to $x$.
Applying these differences to any convergent expansion proves uniqueness.
\end{proof}

\begin{lemma}[Totality and independence]\label{characters}
For every finite positive measure $\nu$ on $\R$,
$\overline{\spn}\{e_q:q\in\Q\}=L^2(\nu)$.
If $\nu\ne0$ is atomless, every finite list of distinct exponentials
is linearly independent in $L^2(\nu)$.
\end{lemma}
\begin{proof}
If $h$ is orthogonal to the rational exponentials, then $h\in L^1(\nu)$
by Cauchy--Schwarz. Dominated convergence makes
$F(t)=\int h(x)e_{-t}(x)\,d\nu(x)$ continuous, so $F=0$ on $\R$.
For $a=\operatorname{Re}h$ its transform is
$(F(t)+\overline{F(-t)})/2=0$, and the same argument applies to
$\operatorname{Im}h$. The finite positive measures $a_+\nu$ and $a_-\nu$
therefore have equal characteristic functions. Mathlib's uniqueness
theorem for finite positive measures makes them equal. On $\{a>0\}$
this implies $\int a_+\,d\nu=0$, and similarly for $a_-$; thus $a=0$.
The imaginary part also vanishes. The orthogonal-complement criterion
for density proves the first assertion. The frequency change $t\mapsto2\pi t$
matches Mathlib's characteristic-function convention.

For distinct $t_1,\ldots,t_d$, let $p(x)=\sum_jc_je_{t_j}(x)$.
If its first $d$ derivatives at any $x$ were zero, the Vandermonde matrix
$((2\pi it_j)^k)_{0\le k<d,1\le j\le d}$ would annihilate
$(c_je_{t_j}(x))_j$. Its determinant is
$\prod_{i<j}(2\pi it_j-2\pi it_i)\ne0$ (factor each difference in the
alternating determinant and compare the leading coefficient).
Thus for nonzero coefficients each zero of $p$ has finite order;
the convergent exponential power series makes that zero isolated.
An isolated subset of $\R$ is countable: assign to each point a rational
interval containing it and no other point of the subset.
An atomless measure gives this zero set measure zero. Since $\nu\ne0$,
$p$ cannot vanish $\nu$-almost everywhere.
\end{proof}

\begin{lemma}[A scalar spectral measure]\label{bochner}
If $\phi:\R\to\C$ is continuous, $\phi(0)=1$, and
$\sum_{j,k}c_j\overline{c_k}\phi(s_j-s_k)\ge0$ for every finite choice
of $c_j,s_j$, then a Borel probability $\mu$ satisfies
$\phi(t)=\int e_t\,d\mu$ for all $t$.
\end{lemma}
\begin{proof}
The two-by-two positive matrices give
$\phi(-t)=\overline{\phi(t)}$ and $|\phi(t)|\le1$.
For $T>0$ put $f_T(t)=(1-|t|/T)_+\phi(t)$ and
$q_T(x)=\int f_T(t)e_{-x}(t)\,dt$. A change of variables gives
\[
 q_T(x)=\frac1T\int_0^T\int_0^T
 \phi(s-t)e^{-2\pi ix(s-t)}\,ds\,dt\ge0.
\]
Indeed the Riemann sums of the continuous integrand are nonnegative
quadratic forms from the hypothesis. For $a>0$, Fubini and the Gaussian
Fourier formula give
\[
 \int q_T(x)e^{-\pi a x^2}\,dx
 =\int f_T(t)a^{-1/2}e^{-\pi t^2/a}\,dt\longrightarrow f_T(0)=1
 \quad(a\downarrow0).
\]
The integrands for Fubini are absolutely integrable because $f_T\in L^1$
and the Gaussian is integrable. The last limit follows after $t=\sqrt a\,u$
by dominated convergence, since $|f_T|\le1$ and the Gaussian has integral one.
Monotone convergence on the left shows $\int q_T=1$.
Fourier inversion, applicable to the continuous integrable $f_T$ with
integrable Fourier transform $q_T$, now gives
$\int e_t(x)q_T(x)\,dx=f_T(t)$.

The probabilities $\mu_n=q_n(x)\,dx$ are tight. To see this, Fubini gives,
with the value at $x=0$ interpreted by continuity,
\[
 \int\left(1-\frac{\sin(2\pi\delta x)}{2\pi\delta x}\right)d\mu_n(x)
 =\frac1\delta\int_0^\delta(1-\operatorname{Re}f_n(t))\,dt.
\]
The integrand on the left is nonnegative and is at least $1/2$ for
$|x|\ge1/\delta$. For all $n\ge1$ and $0\le t\le\delta<1$,
$|1-f_n(t)|\le|1-\phi(t)|+\delta$.
Continuity at zero therefore makes the tail masses uniformly small.
Mathlib's tightness--compactness theorem supplies a cluster probability
$\mu$ and a subnet converging weakly to it. Since $e_t$ is bounded and
continuous and $f_n(t)\to\phi(t)$, weak convergence gives
$\int e_t\,d\mu=\phi(t)$, as required.
\end{proof}

\section{Finite Hilbert-space estimates}
For an ordered independent family $v=(v_1,\ldots,v_d)$, its Gram matrix
is $G_{ij}=\ip{v_j}{v_i}$. Its basis constant is
\[
 K(v)^2=\max_{0\le k\le d}\sup_{c\ne0}
 \frac{c^*D_kGD_kc}{c^*Gc},
\]
where $D_k$ keeps the first $k$ coordinates. Multiplying the measure by
a positive scalar does not change this constant.

\begin{lemma}[Matrix and projection facts]\label{matrixfacts}
The following rules hold.
\begin{enumerate}
\item If $G\le L\le(1+a)G$, then $K(L)\le\sqrt{1+a}K(G)$.
Taking a subsequence in its inherited order cannot increase $K$.
Removing an initial block and orthogonally projecting the remaining
vectors off its span cannot increase $K$ either.
\item If two spans have angle bound
$|\ip{x}{y}|\le\beta\|x\|\|y\|$, with $\beta<1$, concatenating their
bases gives constant at most
$\sqrt{(1+\beta)/(1-\beta)}\max(K_1,K_2)$.
\item For a positive definite matrix $X$,
\[
 v^*X^{-1}v=\sup_z\{2\operatorname{Re}(v^*z)-z^*Xz\}.
 \tag{*}
\]
Thus inversion reverses order; inverse compression is bounded by
compression of the inverse; and
$(\int X\,d\pi)^{-1}\le\int X^{-1}\,d\pi$
for a probability average of positive matrices with a common positive lower bound.
\item Every nonzero bounded projection $P$ on a Hilbert space satisfies
$\|P\|^2=1+\|P-P^*\|^2$; if $P\ne I$ then $\|I-P\|=\|P\|$.
\item Positive definite Gram matrices varying continuously have continuous
basis constants and continuous squared distances from any additional vector
whose inner products also vary continuously.
\end{enumerate}
\end{lemma}
\begin{proof}
The first comparison follows by applying the quadratic-form bounds to
$D_kc$ and $c$. Restricting $c$ to the selected coordinates proves the
subsequence claim. For the initial-span assertion, write a tail vector as
$x$, choose $a$ in the initial span with $\|a+x\|=\dist(x,V)$, and apply
the old prefix bound to $a+x$. Projecting its prefix off $V$ gives exactly
the desired tail prefix; its norm is at most $K\dist(x,V)$.
For (2) use
\[
 (1-\beta)(\|x\|^2+\|y\|^2)\le\|x+y\|^2
 \le(1+\beta)(\|x\|^2+\|y\|^2)
\]
and apply the two prefix bounds in the orthogonal direct-sum norm.
In particular, orthogonal concatenation has constant at most the maximum.
Completing the square proves (*). Restricting the supremum proves the
compression assertion, and moving a supremum outside an average proves
the inverse-average inequality. For (4), relative to
$\operatorname{ran}P\oplus(\operatorname{ran}P)^\perp$, write
$P=\left(\begin{smallmatrix}I&A\\0&0\end{smallmatrix}\right)$.
Its squared norm is $1+\|A\|^2$, its skew part has norm $\|A\|$, and
$I-P$ has the same norm when its range is nonzero.
Finally, near a positive matrix the Rayleigh quotients defining $K$
converge uniformly on the Euclidean unit sphere, because their
denominators have a common positive lower bound. If $g_i=\ip{f}{v_i}$,
the squared distance is $\|f\|^2-g^*G^{-1}g$.
Inverse continuity follows from
$Y^{-1}-X^{-1}=Y^{-1}(X-Y)X^{-1}$ and the same lower bound.
\end{proof}

\begin{lemma}[Conditional means and projections]\label{information}
Let $d\sigma=M^{-1}W(s,y)\,d\rho(s)\,dy$ where $\rho$ is a probability,
$W>0$, and $\int W(s,y)\,dy=M$ does not depend on $s$.
For $u\in L^2(\sigma)$ define its conditional mean explicitly by
$\E[u\mid s]=M^{-1}\int u(s,y)W(s,y)\,dy$; Cauchy--Schwarz in $y$
shows that this belongs to $L^2(\rho)$.
For a finite family $(u_k)$ in $L^2(\sigma)$ let $T$ send its coefficients to its
conditional mean, let $C>0$ be the Gram matrix of its centered parts,
and set $J=TC^{-1}T^*$. For $F=\spn(u_k)$ and $v\in L^2(\rho)$,
\[
 \|P_Fv\|^2=\ip{\Psi(J)v}{v},\qquad \Psi(J)=J(I+J)^{-1}.
\]
For conditionally independent families, their centered Gram matrices
combine diagonally and their operators $J$ add. Furthermore
$\|\Psi(J)-\Psi(J')\|\le\|J-J'\|$ for positive $J,J'$.
\end{lemma}
\begin{proof}
Fubini shows that every centered part is orthogonal to every base function.
Thus the full Gram matrix is $C+T^*T$, and solving the normal equations
for the orthogonal projection gives
$T(C+T^*T)^{-1}T^*$ on the base space.
Put $A=TC^{-1/2}$; multiplying
$(I+AA^*)A=A(I+A^*A)$ by the inverses proves that this is
$AA^*(I+AA^*)^{-1}$. For independent blocks the integral of a product
of centered parts is zero by iterated integration, proving additivity.
Finally $\Psi(J)=I-(I+J)^{-1}$, and the resolvent identity proves the
bound, since both inverse factors have norm at most one.
\end{proof}

\section{A scalar family with small mass loss}
For an even weight $0<w\le1$ bounded away from zero on a torus and
distinct nonconstant characters $\chi_1,\ldots,\chi_d$, write
\[
 B=\begin{pmatrix}m&b^*\\b&G\end{pmatrix},\quad m=\int w=1-\epsilon,
 \quad A=G-bb^*/m,\quad C=A/m,\quad a=b/m.
\]
Here $B$ is the unnormalized Gram matrix of $(1,\chi_1,\ldots,\chi_d)$,
so $0<B\le I$ and every diagonal entry is $m$. Evenness makes $B$ real.
Under $w/m$, a polynomial with coefficient vector $c$ has mean $a^*c$
and variance $c^*Cc$. Define its gain by $Z=a^*C^{-1}a$.
Completing a square, or multiplying the two block matrices, gives
\begin{equation}\label{scalaridentities}
 \min_{a^*c=1}c^*Cc=Z^{-1},\quad
 a^*(C+aa^*)^{-1}a=\frac{Z}{1+Z},\quad m(B^{-1})_{00}=1+Z.
\end{equation}
The first identity is used only when $Z>0$; its minimizer is $C^{-1}a/Z$.

\begin{lemma}[Splitting a row]\label{rowsplit}
Let $B$ be as above and set $H=B^{-1}$.
Put $W(x,y)=1-\mathbf1_D(y)(1-w(x))$, where $D\subset\T$ is symmetric
with Haar measure $p$. Choose any finite list $Y$ of distinct characters
in $y$, starting with $1$. Order the products $\chi_j\psi$ by rows
$j=0,\ldots,d$, where $\chi_0=1$ and $\psi\in Y$.
If $F_j$ is the squared prefix norm through row $j$ in the original
augmented family, and $F_{-1}=1$, then a prefix ending in row $j$
of the product family has squared norm at most
\[
 \max(F_{j-1},F_j)+\frac{\epsilon}{4}(H_{jj}-1).
\]
This is also a bound for its centered family.
\end{lemma}
\begin{proof}
On $\C^{d+1}\otimes L^2(\T)$ the Gram operator and its inverse are
\begin{equation}\label{rareoperators}
 \mathcal G=I+(B-I)\otimes M_D,\qquad
 \mathcal G^{-1}=I+(H-I)\otimes M_D,
\end{equation}
where $M_D$ is multiplication by $\mathbf1_D$.
On this full space consider the operator that is the identity in
completed rows, keeps the selected characters in the current row, and
is zero in later rows. Its restriction to the finite product family is
the prefix under consideration. For the current row let $\Pi$ be the
orthogonal projection onto its selected finite set of characters.
Diagonalize $\Pi M_D\Pi$ on $\operatorname{ran}\Pi$.
If $v$ is a unit eigenvector with eigenvalue $0<t<1$, then
$z=(M_Dv-tv)/\sqrt{t(1-t)}$ is a unit vector in $\ker\Pi$.
The vectors so obtained are orthogonal for different $v$, and on their
two-dimensional spans
\[
 \Pi=\begin{pmatrix}1&0\\0&0\end{pmatrix},\qquad
 M_D=\begin{pmatrix}t&\sqrt{t(1-t)}\\\sqrt{t(1-t)}&1-t\end{pmatrix}.
\]
Eigenvalues zero and one give common eigenspaces. On the orthogonal
remainder $\Pi=0$ and $M_D$ is a projection. This proves the needed
two-projection decomposition directly.

Let $P$ be the set of rows preceding $j$. After discarding identity
summands, the Gram matrices of the range vectors and their dual vectors
on such a block are respectively
\[
 C_t=\begin{pmatrix}B_{PP}&\sqrt t B_{Pj}\\
 \sqrt t B_{jP}&1-t+tm\end{pmatrix},\quad
 D_t=\begin{pmatrix}H_{PP}&\sqrt t H_{Pj}\\
 \sqrt t H_{jP}&1-t+tH_{jj}\end{pmatrix}.
\]
The squared norm of a coordinate projection with these two Gram matrices
is $\|C_t^{1/2}D_tC_t^{1/2}\|$: writing the projection as range synthesis
followed by dual analysis proves this by multiplication.
Here is a bound that avoids any interpolation theorem.
Put $x=B^{1/2}e_j$, $y=B^{-1/2}e_j$,
$P_0=B^{1/2}E_PB^{-1/2}$, $P_1=P_0+xy^T$ and $s=\sqrt{t(1-t)}$.
The projection
\[
 R_t=\begin{pmatrix}P_0+txy^T&sx\\sy^T&1-t\end{pmatrix}
\]
has the preceding range and dual Gram matrices: use the vectors
$(B^{1/2}e_k,0)$ for $k\in P$ and $(\sqrt t x,\sqrt{1-t})$,
and the analogous vectors with $B^{-1/2}$. Define
$S_i=P_i-P_i^T$, $v=x-y$ and $a_0^2=mH_{jj}-1\ge0$.
Then $A_0=S_1-S_0=xy^T-yx^T$ and
\[
 R_t-R_t^T=\begin{pmatrix}(1-t)S_0+tS_1&sv\\-sv^T&0\end{pmatrix},
 \quad \|v\|^2-a_0^2=(1-m)(H_{jj}-1).
\]
For a real skew matrix $S$ and real $b$, one has
\begin{equation}\label{arrow}
 \left\|\begin{pmatrix}S&sb\\-sb^T&0\end{pmatrix}\right\|^2
 =\lambda_{\max}(-S^2+s^2bb^T).
\end{equation}
Indeed, for $q>\|S\|^2$ the Schur complement of
$\sqrt q I-iS$ says that the norm on the left before squaring is
at most $\sqrt q$ exactly when
$s^2b^T(qI+S^2)^{-1}b\le1$; the imaginary skew part has zero quadratic
form on real $b$. The rank-one Schur complement makes this equivalent
to $qI+S^2-s^2bb^T\ge0$. Both sides of (\ref{arrow}) are at least
$\|S\|^2$, which proves the identity by varying $q$.
Consequently increasing $s^2$ by $h$ increases this squared norm by
at most $h\|b\|^2$.

If $a_0>0$, then $A_0^2=-a_0^2I$ on $\spn(x,y)$.
The unit vector $u=-A_0v/(a_0\|v\|)$ satisfies
$A_0u=a_0v/\|v\|$.
Conjugate $\diag(S_0,S_1)$ by the orthogonal block rotation with rows
$(\sqrt{1-t}I,\sqrt t I)$ and $(-\sqrt t I,\sqrt{1-t}I)$,
then compress its second block to $\C u$.
The resulting skew matrix has upper-left block $(1-t)S_0+tS_1$ and
last column $sa_0v/\|v\|$, and its norm is at most
$\max(\|S_0\|,\|S_1\|)$.
Formula~(\ref{arrow}) bounds the cost of replacing this column by $sv$
by $t(1-t)(1-m)(H_{jj}-1)$. If $a_0=0$, then $A_0=0$, and the same
conclusion follows from (\ref{arrow}) starting with last column zero.
Apply Lemma~\ref{matrixfacts}(4) to $R_t$ and the nonzero $P_i$.
If $P_0=0$, use the upper bound $1+\|S_0\|^2=F_{-1}=1$.
This proves the assertion since $t(1-t)\le1/4$.
Restriction to the finite product family and then removal of its
initial constant by Lemma~\ref{matrixfacts}(1) finishes the proof.
\end{proof}

\begin{lemma}[Joining copies]\label{joincopies}
Let $V_0,H_1,\ldots,H_L$ be orthogonal Hilbert spaces, and let the
ordered bases of $H_i$ have squared constants at most $Q$.
Replace each basis vector $x\in H_i$ by $x+T_ix$, where $T_i:H_i\to V_0$.
Suppose $\|\sum_iT_iT_i^*\|\le Z$, $\|T_iT_i^*\|\le\eta$, and $Q>1+Z$.
Prefixes of the resulting list, preceded by all of $V_0$, have squared
norm at most $Q+Q^2\eta/(Q-1-Z)$.
\end{lemma}
\begin{proof}
For a prefix inside $H_i$ let $P$ be its old prefix map.
Completed $H_j$ split off as identity summands. On the other components
the new prefix has block matrix
$\left(\begin{smallmatrix}A&-T_i(I-P)\\0&P\end{smallmatrix}\right)$,
where $AA^*=I+\sum_{j>i}T_jT_j^*$.
Thus $\|A\|^2\le1+Z$, $\|P\|^2\le Q$ and
$\|T_i(I-P)\|^2\le Q\eta$ (also for $P=0,I$).
The inequality $\|a-b\|^2\le(1+h)\|a\|^2+(1+h^{-1})\|b\|^2$,
with $h=(Q-1-Z)/(1+Z)$, gives the claimed bound on a vector $(x,y)$.
\end{proof}

\begin{proposition}[Scalar construction]\label{scalar}
For $0<\delta\le1$ there is $c_\delta>0$ such that, for every
$0<z<c_\delta$ and $\eta>0$, an even trigonometric polynomial $0<w\le1$
on a finite torus and distinct ordered nonconstant characters satisfy
\[
 0<\epsilon=1-\int w<\eta,\quad Z=z,\quad
 K(1,\chi)<1+\delta,\quad K(\chi-\E\chi)<1+\delta,
\]
\begin{equation}\label{scalarextra}
 B\ge\epsilon^2I,\qquad (B^{-1})_{kk}\le1+2\epsilon\quad(k\ne0).
\end{equation}
\end{proposition}
\begin{proof}
We first allow bounded measurable weights bounded away from zero.
The following consequences of the row construction will be used.
Its mass is $M=1-p\epsilon$. Inverse compression in
Lemma~\ref{matrixfacts}(3), followed by (\ref{rareoperators}), gives,
for any nonconstant product in row $j$,
\begin{equation}\label{rarediagonal}
 (A_W^{-1})_{kk}\le1+p(H_{jj}-1),\qquad
 \|A_W^{-1}\|\le\|B^{-1}\|,
\end{equation}
where $A_W$ is the unnormalized centered Gram matrix.
The inverse of a centered Gram matrix is the nonconstant principal block
of the inverse augmented Gram matrix, by block multiplication.
As $Y$ exhausts all characters of $\T$, the gain increases to
\begin{equation}\label{raregain}
 z_p=M[1+p(H_{00}-1)]-1
 =\frac{p}{1-\epsilon}\{Z+\epsilon^2-p\epsilon(Z+\epsilon)\}.
\end{equation}
To justify the limit, the finite product spaces are dense in the Hilbert
space of (\ref{rareoperators}), because torus characters are total and
the new norm is equivalent to the reference norm. The variational
formula (*) makes the $(0,0)$ entries of their inverse compressions
increase to the full inverse entry. Then use (\ref{scalaridentities}).

Set $a=\delta/400$. The initial weight and single character are
$1-a+a\cos(2\pi x)$ and $e^{2\pi ix}$. Their parameters satisfy
\[
 \epsilon_0=a,\quad B_0=\begin{pmatrix}1-a&a/2\\a/2&1-a\end{pmatrix},
 \quad Z_0=\frac{a^2}{4-8a+3a^2},\quad a^2/4\le Z_0\le a^2/3.
\]
Direct inversion gives (\ref{scalarextra}), and the centered family has
constant one. Put $Q_0=1+\delta/4$ and $\eta_0=Z_0$.
We inductively maintain
\begin{gather*}
 0\le Z_0-Z_n\le Z_0(1-2^{-n})/4,\quad
 1+\delta/4\le Q_n<1+\delta/2,\quad K_{\rm cent}^2\le Q_n,\\
 0<\epsilon_n\le a,\quad B_n\ge\epsilon_n^2I,\quad
 (B_n^{-1})_{kk}\le1+2\epsilon_n\ (k\ne0).
\end{gather*}
For $n\ge1$ the family is a concatenation of independent copies,
each with gain at most $\eta_n$; the initial family is one such block.
For brevity write $\epsilon,Z,B$ at a given stage, and
$C_0=\max_j((B^{-1})_{jj}-1)$, including $j=0$.
Choose $p>0$ so small that
\[
 p\le1/100,\quad p\le\epsilon/(Z+\epsilon),\quad
 pC_0\le\epsilon/100,\quad z_p\le Z_0\cdot2^{-n-3}.
\]
Take a finite $Y$ with gain $z$ between $(1-\epsilon/4)z_p$ and $z_p$,
which is possible by (\ref{raregain}). Make $L=\lfloor Z/z\rfloor$
independent copies of this row construction in disjoint variables,
and concatenate their nonconstant characters. Then
\begin{equation}\label{scalarrecurrence}
 1-\epsilon'=M^L,\qquad Z'=Lz,\qquad \eta_{n+1}=z.
\end{equation}
The identities for $Z'$ and the centered Gram matrix follow by integrating
over the independent variables. In particular $0\le Z-Z'<z$.
The choice of $p$ and (\ref{raregain}) give
\[
 \frac{pZ}{1-\epsilon}\le z_p\le\frac{p(Z+\epsilon^2)}{1-\epsilon},
 \qquad
 \frac1{10}<\frac{(1-\epsilon)Z}{Z+\epsilon^2}-p
 \le Lp\le\frac{1-\epsilon}{1-\epsilon/4}\le1-3\epsilon/4.
\]
For the numerical lower bound use $Z\ge3a^2/16$, $\epsilon\le a\le1/400$
and $p\le1/100$. The elementary bounds
$Lt/(1+Lt)\le1-(1-t)^L\le Lt$, applied with $t=p\epsilon$, yield
\begin{equation}\label{epsrecurrence}
 \epsilon/12\le\epsilon'\le\epsilon-3\epsilon^2/4.
\end{equation}
The first bound follows from $(1-t)^{-L}\ge1+Lt$; the second is
Bernoulli's inequality. Consequently $\epsilon_n\downarrow0$ and
$\sum_n\epsilon_n^2\le4a/3$.

For completeness we check both inverse bounds, which are needed later.
For each nonconstant coordinate of the new augmented matrix $B'$, block
inversion and (\ref{rarediagonal}) give
\[
 ((B')^{-1})_{kk}\le\frac{M(1+pC_0)}{1-\epsilon'},\qquad
 ((B')^{-1})_{kk}-1
 \le\frac{\epsilon'+\epsilon/100}{9/10}<2\epsilon'.
\]
Write $h=\|B^{-1}\|\le\epsilon^{-2}$.
The new normalized centered inverse is
\[
 D_0=M\diag(A_W^{-1},\ldots,A_W^{-1})\le MhI.
\]
If $\mathbf a$ is the new mean vector, then
$\mathbf a^*D_0\mathbf a=Z'\le Z_0$ and
\[
 (1-\epsilon')(B')^{-1}
 =\begin{pmatrix}1+Z'&-\mathbf a^*D_0\\-D_0\mathbf a&D_0\end{pmatrix}.
\]
Cauchy--Schwarz in the $D_0$ metric bounds its quadratic form by the
two-by-two matrix with diagonal $1+Z',Mh$ and off-diagonal
$\sqrt{MhZ'}$. Its largest eigenvalue is at most its trace, so
$\|(B')^{-1}\|\le(h+1+Z_0)/(1-\epsilon')$.
For $\epsilon'\le\epsilon(1-3\epsilon/4)$ a direct subtraction gives
\[
 \frac{1-\epsilon'}{(\epsilon')^2}-\frac1{\epsilon^2}
 \ge\frac{1/2+3\epsilon/16}{\epsilon(1-3\epsilon/4)^2}
 \ge\frac1{2\epsilon}>1+Z_0.
\]
Thus $B'\ge(\epsilon')^2I$.

To bound prefixes, apply Lemma~\ref{joincopies} to the old centered
blocks and their means in the constant functions. Since
$Q_n^2/(Q_n-1-Z_n)<20/\delta$, the old augmented constant squared is
at most $Q_n+20\eta_n/\delta$.
Lemma~\ref{rowsplit} adds at most $\epsilon_n^2/2$ for a nonconstant row.
For the constant row its bound is
\[
 1+Z+\frac{\epsilon(Z+\epsilon)}{4(1-\epsilon)}
 \le1+a^2<Q_n.
\]
Centering and orthogonal concatenation of the new centered copies give
$Q_{n+1}=Q_n+20\eta_n/\delta+\epsilon_n^2/2$.
Since $\sum_n\eta_n<3Z_0/2$, this recursion satisfies
\[
 Q_n\le1+\delta/4+30Z_0/\delta+2a/3
 \le1+\delta/4+\delta/16000+\delta/600<1+\delta/2.
\]
The augmented squared constants are also less than
$1+\delta/2+\delta/24000<(1+\delta)^2$.
This completes the induction, with $Z_n\ge3Z_0/4$ and mass loss tending to zero.

It remains to obtain polynomial weights and the exact prescribed gain.
For a fixed constructed weight convolve with the product Fej\'er kernel
\[
 K_N(y)=\prod_{j=1}^r\frac1{N+1}
 \left|\sum_{k=0}^N e^{2\pi iky_j}\right|^2.
\]
It is nonnegative, even, has integral one, and has finite Fourier support.
Thus $w*K_N$ is an even trigonometric polynomial with the same mass
and with $0<w*K_N\le1$. Its Gram matrix is
$B_N=\int U(y)^*BU(y)K_N(y)\,dy$, where $U$ is the diagonal matrix of
the listed characters (inversion of $y$ has no effect because $K_N$ is even).
This averaging preserves $D_kBD_k\le K^2B$ and $B\ge\epsilon^2I$.
Lemma~\ref{matrixfacts}(3) bounds
$B_N^{-1}\le\int U(y)^*B^{-1}U(y)K_N(y)\,dy$,
so its inverse diagonal does not increase.
The multiplier on each fixed Fourier coefficient tends to one, whence
$B_N\to B$ and $Z(B_N)\to Z(B)\ge3Z_0/4$.
Choose $N$ with $Z(B_N)>Z_0/2$ and set $c_\delta=Z_0/2$.
Finally replace the smoothed weight by $w_\theta=1-\theta+\theta w$,
$0\le\theta\le1$. Its Gram matrix is $(1-\theta)I+\theta B$,
which preserves all prefix bounds and is at least $(\theta\epsilon)^2I$.
The inverse-average inequality gives
$(B_\theta^{-1})_{kk}\le1+2\theta\epsilon$ for $k\ne0$.
Gain is continuous in $\theta$, is zero at zero and exceeds $c_\delta$
at one. The intermediate value theorem supplies gain $z$.
The new mass loss $\theta\epsilon$ is positive and no larger than the
original arbitrarily small loss. This proves the proposition.
\end{proof}

\section{Selecting a prescribed direction}

\begin{lemma}[A choice of scale]\label{scale}
Suppose $R\ge2$, $0<\epsilon<1$, and
$\log(1/\epsilon)\ge8\log(4R)$.
Let $\Phi:[0,1]\to[0,\infty)$ be decreasing, with
\[
 \Phi(\epsilon^2)>\epsilon^{1/4},\quad
 \Phi(\sqrt\epsilon)\le\sqrt\epsilon,\quad
 \Phi(R\beta)\ge\Phi(\beta)/(4R)\quad(R\beta\le1/2).
\]
There is $m\in[\epsilon^2R,R\sqrt\epsilon]$ such that
$\sqrt\epsilon<\Phi(m/R)\le\epsilon^{1/4}$ and
$\Phi(m)\le R^{-1/16}\Phi(m/R)$.
\end{lemma}
\begin{proof}
Use the grid $x_k=\epsilon^2R^k$, stopping at the first $K$ with
$x_K\ge\sqrt\epsilon$. Then $x_K\le R\sqrt\epsilon\le1/2$ and
$K\le3\log(1/\epsilon)/(2\log R)+1$.
Let $a,b$ be the first crossings of the levels $\epsilon^{1/4}$ and
$\epsilon^{1/2}$. The ratio bound gives
$\Phi(x_a)>\epsilon^{1/4}/(4R)\ge\epsilon^{1/2}$, so $a<b$.
If every step between them lost a factor smaller than $R^{1/16}$, then
\[
 \tfrac14\log(1/\epsilon)-\log(4R)
 <\log\frac{\Phi(x_a)}{\Phi(x_b)}
 <\tfrac3{32}\log(1/\epsilon)+\tfrac1{16}\log R,
\]
contrary to the assumed logarithmic inequality.
At a step $k$ where this fails, $m=x_k$ has the stated properties.
\end{proof}

\begin{proposition}[Directional construction]\label{selective}
Let $\rho$ be a probability on $\R\times\T^r$, and let
$V=\spn(\phi_1,\ldots,\phi_d)$ have an ordered basis of characters
with constant $K_V$. Fix a unit vector $e\in V$, $K_s>1$ and $\sigma>0$.
For every $\eta>0$ there are a trigonometric polynomial $0<W\le1$,
with constant auxiliary integral $M>1-\eta$, and a finite ordered family
of characters with nonzero new auxiliary frequencies, such that its
conditional means lie in $V$ and
\[
 \|J-\sigma P_{\C e}\|<\eta,\qquad
 K_{\rm cent}^2\le
 1+\bigl(\sqrt{K_s^2-1}+\sqrt{K_V^2-1}\bigr)^2+\eta.
\]
The real frequencies of the new characters are inherited from the
$\phi_j$ or equal zero. If the $\phi_j$ have pairwise distinct old
auxiliary frequencies, the new family has pairwise distinct full
auxiliary frequencies, including all old and new variables.
\end{proposition}
\begin{proof}
Choose $0<\delta\le1$ with $1+\delta<K_s$, and fix $0<z<c_\delta$.
Proposition~\ref{scalar} gives a sequence of scalar systems with this
gain, whose mass losses $\epsilon$ tend to zero. For their blocks $b,G$ put
\[
 \Phi(\beta)=(1-\beta)^2b^*[\beta I+(1-\beta)G]^{-1}b.
\]
This is the unnormalized squared projection of $1$ onto the scalar
nonconstant family for weight $\beta+(1-\beta)w$.
Diagonalizing $G$ shows that $\Phi$ is decreasing: each summand is
a nonnegative multiple of $(1-\beta)^2/(\beta+(1-\beta)g)$, with
$0<g\le1$. For small $\epsilon$,
\begin{equation}\label{responsebounds}
 \Phi(\epsilon^2)\ge\frac{z}{8(1+z)},\qquad
 \Phi(\sqrt\epsilon)\le\sqrt\epsilon,\qquad
 \Phi(R\beta)\ge\frac{\Phi(\beta)}{4R}quad(R\beta\le1/2).
\end{equation}
Indeed $\Phi(0)=(1-\epsilon)z/(1+z)$,
$G\ge\epsilon^2I$, and
$\epsilon^2I+(1-\epsilon^2)G\le2G$.
Also $b_k=-\int\overline{\chi_k}(1-w)$, so the finite orthonormal
projection inequality in Haar $L^2$ gives $\|b\|^2\le\int(1-w)^2\le\epsilon$.
The last inequality follows from
$R\beta I+(1-R\beta)G\le R[\beta I+(1-\beta)G]$ and
$(1-R\beta)^2/(1-\beta)^2\ge1/4$.

Write $e=\sum c_j\phi_j$ and choose
$\kappa=[4\max(1,\sum|c_j|)]^{-1}$; this bounds $|\kappa e(s)|$ by
$1/4$ at every base point, including outside the support of $\rho$.
Set $\theta=\kappa^2$, $a=1-\theta$ and, for $m\le2^{-L}$, define
\[
 w_e(s,y)=m\prod_{\ell=1}^L
 [1+2\kappa\operatorname{Re}(e(s)e^{-2\pi iy_\ell})],
 \qquad \zeta_\ell=e^{2\pi iy_\ell}.
\]
Every factor lies in $[1/2,3/2]$, so $0<w_e\le1$ and
$\int w_e(s,y)\,dy=m$.
Under its conditional probability the means of $\zeta_\ell$ are
$\kappa e(s)$ and the mixed moments for different indices are
$\theta|e(s)|^2$, by direct integration.
In the reference space $L^2(\rho\,dy)$ compress multiplication by $w_e$
to $V\oplus\spn(\zeta_1,\ldots,\zeta_L)$, and call the resulting
positive operator $B_e$. It is $mI$ on $V\cap e^\perp$, $maI$ on
the differences of the $\zeta_\ell$, and on the orthonormal pair
$e,L^{-1/2}\sum_\ell\zeta_\ell$ its matrix is
\[
 mH_L=m\begin{pmatrix}1&\sqrt{L\theta}\\
 \sqrt{L\theta}&a+L\theta\end{pmatrix}.
\]
Its eigenvalues $\gamma<1<\Gamma$ satisfy $\gamma\Gamma=a$ and
$\gamma+\Gamma=1+a+L\theta$. Thus $R=1/\gamma$ is comparable to $L$,
and the squared component of $e$ in the $\gamma$ eigenspace is
$p_\gamma=(\Gamma-1)/(\Gamma-\gamma)\to1$.

Take $L=\lfloor\log(1/\epsilon)/(4\log2)\rfloor$.
For all sufficiently small $\epsilon$, (\ref{responsebounds}) and
Lemma~\ref{scale} give a choice $m\le R\sqrt\epsilon$ with
\begin{equation}\label{window}
 \sqrt\epsilon<\Phi_*:=\Phi(m\gamma)\le\epsilon^{1/4},\qquad
 \Phi(m)/\Phi_*\le R^{-1/16}\longrightarrow0.
\end{equation}
Here $R\asymp L$ verifies the logarithmic condition, and
$m\le R\sqrt\epsilon=o(\epsilon^{1/4})\le2^{-L}$ as required.

Use disjoint variables for $w$ and $w_e$, and put
$W=w+(1-w)w_e$. Its auxiliary integral is
$M=1-(1-m)\epsilon$, independent of $s$.
The new family consists, column by column in $k$, of
\[
 \phi_1\chi_k,\ldots,\phi_d\chi_k,
 \zeta_1\chi_k,\ldots,\zeta_L\chi_k.
\]
Its conditional means are $(1-m)b_k\phi_j/M$ and
$-m\kappa b_ke/M$, respectively; the $b_k$ are real by evenness of $w$.
To compute its projection on $V$, diagonalize $B_e$ in the reference
inner product. On an eigenvector with eigenvalue $\beta$, the new
scalar Gram block is $\beta I+(1-\beta)G$ and its pairing with a
vector $v\in V$ is $(1-\beta)b$ times the component of $v$ on that
eigenvector. These statements follow by expanding
$W=w_e+(1-w_e)w$ and integrating the disjoint variables.
Therefore the unnormalized squared projection form on $V$ is
$P_V\Phi(B_e)P_V$. It equals $\Phi(m)$ on $V\cap e^\perp$ and equals
\[
 \Phi_e=p_\gamma\Phi(m\gamma)+(1-p_\gamma)\Phi(m\Gamma)
 \quad\text{on }\C e.
\]
For small $\epsilon$, $3\Phi_*/4\le\Phi_e\le\Phi_*$.
The centered Gram matrix is positive definite: a zero centered
combination would, since $W>0$, be independent of all auxiliary
variables under the reference measure. Its nonconstant $\chi_k$
Fourier coefficients would then vanish. Independence of $\phi_j$ in
$L^2(\rho)$ and of the separate $\zeta_\ell$ forces all coefficients
to be zero. Lemma~\ref{information} now implies
\[
 J=\lambda_\perp P_{V\cap e^\perp}+\lambda P_{\C e},\qquad
 \lambda_\perp=\frac{\Phi(m)}{M-\Phi(m)},\quad
 \lambda=\frac{\Phi_e}{M-\Phi_e}.
\]
The denominators are positive; equivalently the projection form is
strictly below the identity because the centered Gram matrix is positive.
By (\ref{window}),
\begin{equation}\label{selectivelimits}
 \lambda\to0,\qquad \lambda_\perp/\lambda\to0,\qquad
 (1-M)/\lambda\to0.
\end{equation}

We must control prefixes inside each column, not just entire columns.
Temporarily adjoin the constant scalar column with all rows.
Call its Gram matrix $H$ and the reference row Gram matrix
$\Omega=\diag(S_V,I_L)$.
On each reference eigenspace of $B_e$, the full scalar matrix is
$\beta I+(1-\beta)B$, whose prefix maps have norm at most $K_s$.
Hence completed-column prefixes have that bound.
For a nonconstant column $k$, let $I_k$ insert its coefficient vector.
The weight bound $W\le1$, the inverse-average inequality, and
(\ref{scalarextra}), in those same orthonormal eigencoordinates, give
\[
 I_k^*HI_k\le\Omega,\qquad
 I_k^*H^{-1}I_k\le(1+2\epsilon)\Omega^{-1}.
\]
Undoing the change to orthonormal row coordinates gives exactly these
two coefficient-matrix inequalities. Synthesis from this column has
norm at most one in the $\Omega$ metric, while coefficient extraction
into it has norm at most $\sqrt{1+2\epsilon}$.
An initial row projection in that metric has norm at most $K_V$.
Thus a prefix ending in this column is $P+Q$, where $PQ=QP=0$,
$\|P\|\le K_s$ and $\|Q\|\le\sqrt{1+2\epsilon}K_V$.
Apply Lemma~\ref{matrixfacts}(4) and the triangle inequality to its
skew part (with zero projections treated separately) to obtain
\begin{equation}\label{selectedprefix}
 \|P+Q\|^2\le
 1+\bigl(\sqrt{K_s^2-1}+
 \sqrt{(1+2\epsilon)K_V^2-1}\bigr)^2.
\end{equation}
Remove the $\zeta_\ell$ from the constant column by restricting to the
remaining coordinates. All prefixes ending in a nonconstant column
retain this bound. The remaining initial column spans $V$;
orthogonally projecting it away replaces each new vector by its
centered part, since its mean lies in $V$. The quotient argument in
Lemma~\ref{matrixfacts}(1), applied only to these prefixes, proves
(\ref{selectedprefix}) for the centered family.

Finally use $n=\lceil\sigma/\lambda\rceil$ conditionally independent
copies. Their operators add, their centered parts are orthogonal, and
their mass is $M^n$. Formula~(\ref{selectivelimits}) gives
$n\lambda\to\sigma$, $n\lambda_\perp\to0$ and
$1-M^n\le n(1-M)\to0$.
The centered bound (\ref{selectedprefix}) tends to the asserted bound.
Taking $\epsilon$ sufficiently small proves both conclusions.
Every character used has a nonconstant scalar frequency in its copy's
fresh variable block. Helper rows inherit the old frequencies, phase
rows have their separate $y_\ell$ frequencies, and different copies
use disjoint scalar blocks. This also proves the assertions about
distinct auxiliary modes and real frequencies.
\end{proof}

\section{Extending a finite exponential family}
Let $\rho$ be an atomless probability on $\R$. Here and below atomless
means that every singleton has measure zero.
Let $A$ be a nonempty ordered list of distinct rational exponentials
and let $f$ be another such exponential. Put $E=\spn A$ and
\[
 r=f-P_Ef,\quad e=r/\|r\|_\rho,\quad
 0<u\le1/100,\quad K_*=1+10u,\quad c=u/(1+2u).
\]
Lemma~\ref{characters} ensures $r\ne0$.
Write $d_\rho(f,A)=\dist_{L^2(\rho)}(f,E)$.

\begin{lemma}[Copies of the residual]\label{copies}
For every $a>0$ there are a trigonometric polynomial $0<Q\le1$ on a
finite torus, with $\int Q>1-a$, and finite ordered character families
$V_i$ with unit vectors $h_i\in\spn V_i$ in
$L^2(d\rho\,Q\,dy/\int Q)$ such that
\[
 K(V_i)^2\le(1+u)^3,\quad J_0E=0,\quad
 \|J_0\|<2u,\quad \ip{\Psi(J_0)e}{e}\ge c,
 \qquad J_0=u\sum_{i=1}^NP_{\C h_i}.
\]
Each $V_i$ has distinct nonzero auxiliary modes and rational real frequencies.
\end{lemma}
\begin{proof}
Write $\phi=(A,f)$, $r=\sum_j a_j\phi_j$, and let $S$ be its Gram
matrix in $L^2(\rho)$. Choose $0<z<c_u$ with $z\|S\|<u$.
For each $j$ use Proposition~\ref{scalar}, with parameters $u,z$, on
separate torus variables. Denote its nonconstant characters by $\psi_{jk}$.
In the ordered family $(\phi_j\psi_{jk})_{j,k}$ the conditionally centered
Gram matrix is $C=\diag(C_j)$, and its conditional mean map satisfies
\[
 TC^{-1}T^*=z\sum_j\phi_j\otimes\phi_j,
 \quad\|TC^{-1}T^*\|=z\|S\|<u,
 \quad C\le C+T^*T\le(1+u)C.
\]
Here $(v\otimes v)x=\ip{x}{v}v$; the nonzero eigenvalues of the
synthesis operator times its adjoint equal those of its Gram matrix,
as follows by mapping each nonzero eigenvector through synthesis or its
adjoint. Each centered block has constant less than $1+u$; orthogonal
concatenation and Lemma~\ref{matrixfacts}(1) give squared constant
at most $(1+u)^3$ for the uncentered family.

By (\ref{scalaridentities}), choose $g_j\in\spn_k\psi_{jk}$ of mean one
and variance $1/z$. Conditional independence and $|\phi_j|=1$ give
\[
 h=\sum_j a_j\phi_jg_j,\quad \E[h\mid x]=r,\quad
 \|h\|^2=\|r\|_\rho^2+z^{-1}\sum_j|a_j|^2.
\]
Thus $\alpha=\|r\|_\rho/\|h\|$ lies in $(0,1)$ and depends only
on $r,z$, not on the scalar mass losses.
Set $N=\lceil\alpha^{-2}\rceil$ and take $N$ copies of the whole
construction, independent given $x$, normalizing $h$ in each copy.
Their means are $\alpha e$; their centered parts are mutually orthogonal,
so $h_i\perp E$ and $\ip{h_i}{h_j}=\alpha^2$ for $i\ne j$.
Their Gram matrix has largest eigenvalue $1+(N-1)\alpha^2$, whence
\[
 \|J_0\|=u[1+(N-1)\alpha^2]<2u,\qquad
 \ip{J_0e}{e}=uN\alpha^2\ge u.
\]
Diagonalizing the finite-rank positive $J_0$ shows
$\Psi(J_0)\ge J_0/(1+\|J_0\|)$, proving the stated lower bound.
Choose all scalar losses so that their finite product of masses exceeds
$1-a$, and let $Q$ be their product weight. This choice does not alter
$\alpha$ or $N$. Separate variable blocks for different $j$ and copies,
and nonzero modes of every $\psi_{jk}$, give the frequency assertions.
\end{proof}

\begin{lemma}[Product extension]\label{product}
For $\delta,\tau>0$ there are a trigonometric polynomial $0<W\le1$
on $\R\times\T^D$ and an ordered list
$B=(e_{b_i}(x)e^{2\pi ik_i\cdot y})_i$, with $b_i\in\Q$ and distinct
$k_i\in\Z^D\setminus\{0\}$, such that
\[
 \int W(x,y)\,dy=m>1-\delta,\quad d\sigma=m^{-1}W\,d\rho\,dy,
\]
\[
 K_\sigma(A,B)<\max\{K_\rho(A),K_*\}+\tau,\qquad
 d_\sigma(f,(A,B))^2<(1-c/2)d_\rho(f,A)^2.
\]
\end{lemma}
\begin{proof}
Apply Lemma~\ref{copies} with $a=\delta/2$. On the resulting base
probability apply Proposition~\ref{selective} to each $V_i,h_i$, using
$K_s=1+u$, strength $u$, and tolerance $\eta$, in separate new variables.
Let $W=Q\prod_iW_i$ and concatenate their character lists.
The auxiliary mass is the constant $m=(\int Q)\prod_iM_i>1-\delta$
for sufficiently small $\eta$. The frequency conclusion of
Proposition~\ref{selective} applies because each $V_i$ already has
distinct auxiliary modes. Different applications have different fresh
scalar blocks, so all full $k_i$ are distinct and nonzero.

Lemma~\ref{information} gives $J=\sum_iJ_i$ with
$\|J-J_0\|<N\eta$. The centered families are orthogonal. Since
$(1+u)^2-1\le3u$ and $(1+u)^3-1\le4u$, their concatenation has
\[
 K_{\rm cent}^2\le1+(\sqrt3+2)^2u+\eta<1+14u
\]
for small $\eta$. Also $\|J\|<3u$, and hence
$C\le C+T^*T\le(1+3u)C$ for the concatenated Gram matrices.
Thus its uncentered constant satisfies
$K_\sigma(B)^2<(1+3u)(1+14u)<(1+10u)^2=K_*^2$.

For $F=\spn B$ and $\beta=\|P_F|_E\|$,
Lemma~\ref{information} gives
\[
 \beta^2=\|P_E\Psi(J)|_E\|\longrightarrow0,\qquad
 \|P_Fe\|^2=\ip{\Psi(J)e}{e}>c/2
\]
for small $\eta$. Indeed $J_0E=0$ and the resolvent bound applies.
All old inner products are unchanged because the base marginal remains
the original probability. Lemma~\ref{matrixfacts}(2) bounds the
concatenated basis constant by
$\sqrt{(1+\beta)/(1-\beta)}\max\{K_\rho(A),K_*\}$.
Finally $P_Ef+P_Fr\in E+F$, so
$d_\sigma(f,(A,B))^2\le\|r-P_Fr\|^2
=\|r\|_\rho^2(1-\|P_Fe\|^2)$.
These estimates prove both strict inequalities.
\end{proof}

\section{Transfer and passage to a basis}
\begin{lemma}[Graph averaging]\label{graph}
Let $\mu$ be a finite atomless measure on $\R$. For finitely many sums
$G(x,y)=\sum_k g_k(x)e^{2\pi ik\cdot y}$, with finite support and
$g_k\in L^1(\mu)$, and any finite
$L\subset\Z^D\setminus\{0\}$, there exist $\omega\in\Z^D$ and
integers $j_n\uparrow\infty$ such that
\[
 k\cdot\omega\ne0\ (k\in L),\qquad
 \int G(x,j_n\omega x)\,d\mu\longrightarrow\iint G(x,y)\,dy\,d\mu.
\]
\end{lemma}
\begin{proof}
For $D=0$ there is nothing to prove. Otherwise choose
$\omega=(1,b,\ldots,b^{D-1})$ with integer $b$ avoiding the finitely
many roots of the nonzero polynomials associated to $L$ and all nonzero
modes of the sums. For any such mode put $l=k\cdot\omega\ne0$.
If $a_j=\int e_{jl}(x)g(x)\,d\mu(x)$, then
\[
 \frac1M\sum_{j=1}^M|a_j|^2
 =\iint g(x)\overline{g(y)}D_M(l(x-y))\,d\mu(x)d\mu(y),\qquad
 D_M(t)=\frac1M\sum_{j=1}^Me^{2\pi ijt}.
\]
Fubini is justified by integrability of $|g(x)g(y)|$.
The geometric-sum formula gives $D_M(t)\to0$ unless $t\in\Z$,
and $|D_M(t)|\le1$. For each $x$, the exceptional $y$ belong to
$x+l^{-1}\Z$, a countable null set. Fubini and dominated convergence
therefore give $M^{-1}\sum_{j\le M}|a_j|^2\to0$.
Sum these nonnegative quantities over the finitely many modes and
functions. For each $n$ arbitrarily large $j$ have this sum less than
$1/n$; otherwise its Ces\`aro average could not tend to zero.
Choose such $j_n$ successively. All nonconstant terms tend to zero,
which proves the assertion without a Wiener-lemma citation.
\end{proof}

\begin{lemma}[One stage]\label{stage}
Let $\mu$ be an atomless probability, $d\nu=v\,d\mu$ with $0<v\le1$,
and let $A,f$ be as above. For $\delta,\tau>0$ there are $0<V\le1$
and an extension $A'$ by distinct rational exponentials such that
\[
 \int(1-V)\,d\mu<\delta,\quad
 K_{V\nu}(A')\le\max\{K_\nu(A),K_*\}+\tau,
\]
\[
 d_{V\nu}(f,A')^2\le(1-c/4)d_\nu(f,A)^2.
\]
\end{lemma}
\begin{proof}
The finite measure $\nu$ is nonzero and atomless. Apply
Lemma~\ref{product} to $\rho=\nu/\nu(\R)$ with tolerances
$\delta/2,\tau/2$. Apply Lemma~\ref{graph} to
\[
 G\in\{1-W\}\cup\{vW p\overline q:p,q\in(A,B,f)\},\qquad
 L=\{k_i\}_i\cup\{k_i-k_j:i\ne j\}.
\]
These are finite sums with integrable coefficients: all polynomials
are bounded and $0<v\le1$ on the finite measure space.
Set $V_n(x)=W(x,j_n\omega x)$ and
$t_{i,n}=b_i+j_n k_i\cdot\omega$. The slopes are distinct and nonzero;
for sufficiently large $n$ the $t_{i,n}$ are distinct from each other
and from all old frequencies. They are rational.

The product Gram matrix is positive definite. Indeed $W>0$, and
integration in the torus variables separates each nonzero $k_i$ from
the zero mode and from every other $k_j$; within the zero mode use
Lemma~\ref{characters} for $A,f$.
Graph averaging and Lemma~\ref{matrixfacts}(5) now give
\begin{gather*}
 \int(1-V_n)\,d\mu\to1-m,\quad
 K_{V_n\nu}(A,(e_{t_{i,n}})_i)\to K_\sigma(A,B),\\
 d_{V_n\nu}(f,(A,(e_{t_{i,n}})_i))^2
 \to m\nu(\R)d_\sigma(f,(A,B))^2.
\end{gather*}
The limiting distance is less than
$(1-c/2)d_\nu(f,A)^2$, since $m\le1$ and
$\nu(\R)d_\rho(f,A)^2=d_\nu(f,A)^2>0$.
Choose a sufficiently large $n$ to obtain the stated weaker bound
and the other two estimates, and put $V=V_n$.
\end{proof}

\begin{proposition}[Weighted exponential basis]\label{weighted}
For every atomless probability $\mu$ on $\R$ and $\varepsilon>0$ there
are a weight $0<w\le1$ and distinct $t_n\in\Q$ such that
$(e_{t_n})$ is a Schauder basis of $L^2(w\mu)$ with constant
less than $1+\varepsilon$.
\end{proposition}
\begin{proof}
Take $u=\min(1,\varepsilon)/100$, $A_0=(1)$ and $v_0=1$.
Enumerate $\Q$ and visit the first $k$ elements in round $k$.
For stage $n\ge1$ choose $\delta_n=2^{-n}$ and
$\tau_n=(\varepsilon-10u)2^{-n-1}$.
Apply Lemma~\ref{stage} to $\nu=v_{n-1}\mu$ and the visited exponential,
setting $v_n=v_{n-1}V_n$. If the target is already in the span, set
$A_n=A_{n-1}$ and $V_n=1$. The lists extend in order, and
\[
 K_{v_n\mu}(A_n)\le K_*+\sum_{j\le n}\tau_j
 \le K_\infty:=1+10u+(\varepsilon-10u)/2<1+\varepsilon.
\]
Tonelli gives $\sum_n(1-V_n(x))<\infty$ almost everywhere, since its
integral is at most $\sum_n\delta_n<\infty$.
At such $x$, eventually $V_n(x)\ge1/2$, and
$-\log V_n(x)\le2(1-V_n(x))$. Every factor is positive, so
$v_n(x)\downarrow w(x)>0$. Dominated convergence gives $v_n\to w$
in $L^1(\mu)$. On the null exceptional set set $w=1$.

Enlarging a span and decreasing a weight can only decrease the distance
to the span. If $N_n(q)$ counts visits to $e_q$ through stage $n$, the
stage estimate, including skipped zero-distance stages, implies
\[
 d_{w\mu}(e_q,A_n)^2\le d_{v_n\mu}(e_q,A_n)^2
 \le(1-c/4)^{N_n(q)}d_\mu(e_q,A_0)^2\longrightarrow0.
\]
Thus the union of the lists spans a dense subspace by
Lemma~\ref{characters}. It contains infinitely many distinct characters:
otherwise its finite-dimensional, hence closed, span would be all of
$L^2(w\mu)$, contradicting the arbitrarily long independent lists
given by that lemma. Enumerate this union in its construction order.
For a fixed finite combination $p$ and fixed prefix $P_k$, all its
characters occur in $A_n$ for large $n$. The finite prefix bound and
$L^1$ convergence of the weights yield
$\|P_kp\|_{w\mu}\le K_\infty\|p\|_{w\mu}$, because $p$ and $P_kp$
are bounded functions. The sequence is independent by
Lemma~\ref{characters}; Lemma~\ref{basisfacts} proves it is a Schauder
basis with the required constant.
\end{proof}

\section{Proof of the main theorem}
\begin{proof}[Proof of Theorem~\ref{thm:main}]
First suppose $b_n=U_{t_n}g$ is a Schauder basis.
Its span is contained in the orbit span, so $g$ is cyclic.
Its norms are all $\|g\|$, hence it is weakly null by
Lemma~\ref{basisfacts}. If $a\ne0$ were a common eigenvector,
$U_ta=\lambda_ta$ would have $|\lambda_t|=1$ by unitarity.
Consequently $|\ip{b_n}{a}|=|\ip{g}{a}|$, which must be zero.
Every $b_n$ is then orthogonal to $a$, and density implies $a=0$,
a contradiction.

Conversely, choose a cyclic unit vector $v$; a cyclic vector is nonzero
because $H$ is nonzero. The function $\phi(t)=\ip{U_tv}{v}$ is
continuous and positive definite, since
$\sum c_j\overline{c_k}\phi(s_j-s_k)=\|\sum c_jU_{s_j}v\|^2$.
Lemma~\ref{bochner} gives a probability $\mu$ with
\[
 \ip{U_tv}{U_sv}=\phi(t-s)=\int e_{t-s}\,d\mu
 =\ip{e_t}{e_s}_{L^2(\mu)}.
\]
Thus $Se_t=U_tv$ defines a well-defined linear isometry on the span
of the real exponentials: any zero linear combination has zero image
by this Gram identity. Lemma~\ref{characters} makes its domain dense.
For span elements $F_j\to F$, the sequence $SF_j$ is Cauchy; define
$SF=\lim_jSF_j$. The isometry identity makes the limit independent of
the approximants and preserves linearity and norms.
Its range is closed (limits lift to Cauchy sequences) and contains the
dense orbit span, hence $S$ is onto. For every $t$,
$S(e_tF)=U_tSF$: this holds on the exponential span and extends by
continuity, since multiplication by $e_t$ is an isometry.
If $\mu(\{x\})>0$, the nonzero vector $S\mathbf1_{\{x\}}$ is a common
eigenvector with eigenvalues $e_t(x)$. Hence $\mu$ is atomless.

Apply Proposition~\ref{weighted}. Multiplication by $\sqrt w$ is a
unitary $L^2(w\mu)\to L^2(\mu)$: its inverse is division by
$\sqrt w$, and both norm identities follow directly by integration.
Its composition with $S$ maps $e_{t_n}$ to
$U_{t_n}(S\sqrt w)$. A unitary preserves convergence and uniqueness
of expansions. This gives the desired orbit basis with $g=S\sqrt w$
and pairwise distinct real times.
\end{proof}

\appendix
\section{Dependency audit and Mathlib boundary}\label{audit}

\subsection{External inputs in the original proof}
The following list includes the uncited steps as well as the explicit
citations. It records their replacements, rather than treating a citation
or a theorem name as a proof.
\begin{enumerate}
\item The shrinking-basis theorem from Lindenstrauss--Tzafriri,
continuity of coefficient functionals implicit in the definition of a
Schauder basis, and the final basis criterion are replaced by
Lemma~\ref{basisfacts}. No reflexivity or weak compactness theorem is needed.
\item Bochner's representation theorem is replaced by
Lemma~\ref{bochner}. Its positivity, normalization, integrability and
tightness steps are proved explicitly before invoking the Mathlib
Fourier inversion and probability compactness foundations below.
\item Fourier uniqueness for a complex density, continuity in frequency,
totality of rational exponentials, and independence for atomless measures
are handled in Lemma~\ref{characters}. Only the positive finite-measure
uniqueness theorem is imported from Mathlib; the complex-density reduction
and the countability-of-zeros argument are included.
\item The extension from orbit Gram identities to an onto unitary,
its intertwining identity, exclusion of spectral point masses, weighted
$L^2$ unitaries, and transport of bases are proved in the main proof.
\item The finite Gram comparisons, quotient and concatenation bounds,
Schur complements, inverse compression and inverse averaging,
projection skew-part identity, and Gram continuity needed for transfer
are proved in Lemma~\ref{matrixfacts} and the displayed block calculations.
\item The scalar construction previously cited as Proposition~14.1 of
\cite{FLT} is proved in Proposition~\ref{scalar}. The additional inputs
inside that construction are also resolved: row splitting and the
two-projection decomposition are proved in Lemma~\ref{rowsplit}; its
skew-matrix estimate is (\ref{arrow}); independent-copy prefixes are
bounded in Lemma~\ref{joincopies}; smoothing and tuning are proved at
the end of Proposition~\ref{scalar}. In particular there is no appeal
to Halmos's theorem, the Schur product theorem, or operator Jensen's inequality.
\item The directional construction previously cited as Proposition~15.1
of \cite{FLT} is proved in Proposition~\ref{selective}, including its
scale selection, finite spectral calculation, within-column prefix
bounds, and preservation of the necessary frequency distinctions.
\item Conditional Gram decomposition, the projection identity,
addition under conditional independence, and its perturbation estimate
are proved in Lemma~\ref{information}. Residual replication is proved
in Lemma~\ref{copies}; no part of Section~16.1 of \cite{FLT} is assumed.
\item Wiener's lemma from Katznelson is replaced by the direct product
integral argument in Lemma~\ref{graph}, including the common subsequence
and the choice of an integer direction. No decay of individual Fourier
coefficients of an atomless measure is asserted or needed.
\item Positivity of the transferred Gram matrix, rationality and distinctness
of new frequencies, and normalization of distances are justified in
Lemma~\ref{stage}. Positivity of the infinite product, convergence of
weights, the repeated-target argument, infinitude of the selected list,
and passage of prefix bounds to the limit are proved in
Proposition~\ref{weighted}.
\end{enumerate}

\subsection{Foundations checked in Mathlib}
The source scan used the clean Mathlib checkout at commit
\href{https://github.com/leanprover-community/mathlib4/tree/81a5d257c8e410db227a6665ed08f64fea08e997}{\texttt{81a5d257c8e410db227a6665ed08f64fea08e997}}
(Lean \texttt{v4.32.0}). This records a source-level dependency check,
not a Lean verification of the present paper.
The following are the substantive foundations used without reproducing
their Mathlib proofs. Paths are relative to the Mathlib repository.

\begingroup\small
\begin{enumerate}
\item Completeness of $L^2$, orthogonal projections, the density criterion,
and the representation of continuous Hilbert-space functionals:
\ml{MeasureTheory.Lp.instCompleteSpace} in
\ml{Mathlib/MeasureTheory/Function/LpSpace/Complete.lean};
\ml{Submodule.orthogonalProjectionOnto} in
\ml{Mathlib/Analysis/InnerProductSpace/Projection/Basic.lean};
\ml{Submodule.orthogonal_orthogonal_eq_closure} and
\ml{Submodule.topologicalClosure_eq_top_iff} in
\ml{Mathlib/Analysis/InnerProductSpace/Projection/Submodule.lean};
\ml{InnerProductSpace.toDual} in
\ml{Mathlib/Analysis/InnerProductSpace/Dual.lean}.
The finite orthonormal projection inequality used above follows directly
by expanding the squared norm of the orthogonal residual.
\item The bounded inverse theorem used for the sequence space in
Lemma~\ref{basisfacts}:
\ml{ContinuousLinearEquiv.ofBijective} in
\ml{Mathlib/Analysis/Normed/Operator/Banach.lean}.
Finite-dimensional subspaces are closed by
\ml{Submodule.closed_of_finiteDimensional} in
\ml{Mathlib/Topology/Algebra/Module/FiniteDimension.lean}.
\item Finite-dimensional self-adjoint diagonalization:
\ml{LinearMap.IsSymmetric.eigenvectorBasis} and
\ml{LinearMap.IsSymmetric.apply_eigenvectorBasis} in
\ml{Mathlib/Analysis/InnerProductSpace/Spectrum.lean}.
Positive square roots and inverse norm bounds used here are obtained
by applying the respective scalar functions to this finite diagonalization.
\item Dominated convergence and Fubini:
\ml{MeasureTheory.tendsto_integral_of_dominated_convergence} in
\ml{Mathlib/MeasureTheory/Integral/DominatedConvergence.lean}, and
\ml{MeasureTheory.integral_prod} in
\ml{Mathlib/MeasureTheory/Integral/Prod.lean}.
Monotone convergence and interchange of a nonnegative sum and integral:
\ml{MeasureTheory.lintegral_iSup} and
\ml{MeasureTheory.lintegral_tsum} in
\ml{Mathlib/MeasureTheory/Integral/Lebesgue/Add.lean}.
Every use above specifies an integrable majorant, positivity, or the
finite-measure bound needed for the indicated interchange.
\item Gaussian Fourier transform and Fourier inversion:
\ml{fourier_gaussian_pi} in
\ml{Mathlib/Analysis/SpecialFunctions/Gaussian/FourierTransform.lean}, and
\ml{Continuous.fourierInv_fourier_eq} in
\ml{Mathlib/Analysis/Fourier/Inversion.lean}.
Lemma~\ref{bochner} proves integrability of both transforms before
using inversion; the Fourier normalization here is exactly $2\pi$.
\item Uniqueness for finite positive measures:
\ml{MeasureTheory.Measure.ext_of_charFun} in
\ml{Mathlib/MeasureTheory/Measure/CharacteristicFunction/Basic.lean}.
Its complete, second-countable Borel inner-product-space hypotheses
hold for $\R$; its finite measures here are the positive and negative
parts of integrable real densities.
\item Compactness for tight probabilities:
\ml{isCompact_closure_of_isTightMeasureSet} in
\ml{Mathlib/MeasureTheory/Measure/Prokhorov.lean}.
Cluster points are provided by \ml{IsCompact.exists_clusterPt} in
\ml{Mathlib/Topology/Compactness/Compact.lean}.
Weak convergence against bounded continuous complex functions is
\ml{MeasureTheory.ProbabilityMeasure.tendsto_iff_forall_integral_rclike_tendsto}
in \ml{Mathlib/MeasureTheory/Measure/ProbabilityMeasure.lean}.
The real line has the Hausdorff Borel structure required here, and the
uniform tail estimate in Lemma~\ref{bochner} proves tightness.
\item Totality of torus characters, used only in the rare-set inverse limit:
\ml{AddCircle.span_fourierLp_closure_eq_top} in
\ml{Mathlib/Analysis/Fourier/AddCircle.lean}, with period one and $p=2$.
Only one added torus variable is completed at a time; no unproved
multivariable density theorem is needed.
\end{enumerate}

The source scan also inspected
\ml{Mathlib/Analysis/Normed/Module/Bases.lean}, including
\ml{SchauderBasis.exists_norm_proj_le}. That theorem assumes a
Schauder basis with continuous coordinates; Lemma~\ref{basisfacts}
supplies the bridge from the paper's expansion-only definition instead
of silently assuming continuity. The scalar and directional constructions,
and the representation statement of Lemma~\ref{bochner}, were not located
as ready-to-use theorems in this checkout and are not imported from it.
\endgroup

\subsection{Correspondence with the supplied Lean statement}
In \ml{Showcase.lean}, \ml{IsStronglyContinuousGroup} specifies the
composition law and continuity of each orbit. Since each $U_t$ is
invertible, $U_0U_0=U_0$ implies $U_0=I$, and the law then gives
$U_{-t}=U_t^{-1}$. These are precisely the group properties used above.
The definitions \ml{IsCyclicVector}, \ml{IsCommonEigenvector}, and
\ml{HasSimpleContinuousSpectrum} agree with the definitions in Section~1.
Indexing the basis by $0,1,\ldots$ instead of $1,2,\ldots$ is a shift;
pairwise distinctness is the injectivity of the time sequence.
Lemma~\ref{basisfacts} supplies the continuous biorthogonal coordinates
required by Mathlib's \ml{SchauderBasis}. Mathlib's inner product is
linear in the second variable, so the paper's $\ip{x}{y}$ corresponds
to its inner product with arguments reversed. None of these convention
changes alters the target statement. The supplied Lean file is left
unchanged, including its existing proof placeholder.

\begin{thebibliography}{9}
\raggedright
\bibitem{FLT} D.~Freeman, L.~Liehr and M.~A.~Taylor,
\emph{Structured coordinate systems in $L^p(\R)$},
\href{https://arxiv.org/html/2610.18709v2}{arXiv:2610.18709v2}, 2026.
\end{thebibliography}
\end{document}
